%!TEX root = ../template.tex
%%%%%%%%%%%%%%%%%%%%%%%%%%%%%%%%%%%%%%%%%%%%%%%%%%%%%%%%%%%%%%%%%%%%
%% chapter5.tex
%% NOVA thesis document file
%%
%% Chapter 5: Findings
%%%%%%%%%%%%%%%%%%%%%%%%%%%%%%%%%%%%%%%%%%%%%%%%%%%%%%%%%%%%%%%%%%%%

\typeout{NT FILE chapter5.tex}%

\chapter{Findings}
\label{cha:findings}

This chapter reports the grounded account that emerged from the 14 interviews. It is organised by research question rather than as a catalogue of fourteen disconnected labels. The analysis produced \textbf{14 categories} (CAT1--CAT14) and \textbf{four phenomena} (PH1--PH4). Phenomena are cross-cutting claims that several categories feed; they are not a requirement that every category headline a phenomenon. CAT2, CAT11, and CAT14 answer the research questions but do not feed a phenomenon: they contrast or condition other categories.

Quotations are short excerpts from the coded transcripts, identified by case. Portuguese interviews (S7--S10) are quoted from the English working transcripts used in analysis. Claims stay bounded by memo and case identifiers; student evidence is not merged with lecturer interpretation, and design wishes are not treated as tested interventions.

\section{How to read this chapter}
\label{sec:findings_how}

Each category subsection states (i)~what the category is, (ii)~how it varies, including wave~2 properties, (iii)~at least one short quotation, and (iv)~negative or boundary cases where they exist. Table~\ref{tab:cat_roster} lists all fourteen categories before the detailed sections. Table~\ref{tab:cat_rq_summary} at the end of the chapter repeats the map with research questions, one-line claims, and memo identifiers.

\begin{table}[ht]
\centering
\small
\caption{Overview of the fourteen categories. Detailed evidence follows in the sections below.}
\label{tab:cat_roster}
\begin{tabularx}{\textwidth}{l >{\raggedright\arraybackslash}X}
\toprule
\textbf{ID} & \textbf{What it is} \\
\midrule
CAT1 & Getting started is hard without a demonstration or a path already seen. \\
CAT2 & Too many symbols, arrows, or frameworks create doubt. \\
CAT3 & Teacher examples and walkthroughs unlock the work. \\
CAT4 & When teaching is thin, students invent their own process. \\
CAT5 & Feedback arrives too late for the model still to change easily. \\
CAT6 & Who gets detailed feedback depends on asking and on public versus private format. \\
CAT7 & Checks that catch an invalid model are valued. \\
CAT8 & Tools add cost: options, licences, switches, weak interfaces. \\
CAT9 & Groups can kill a diagram or keep it alive. \\
CAT10 & A real reader makes modelling worth the care. \\
CAT11 & Course load leaves little time for demonstrations and feedback loops. \\
CAT12 & Without a compile result, it is hard to judge ``good enough.'' \\
CAT13 & AI is used to explain or compare, not to generate graded models. \\
CAT14 & Moving from one kind of diagram to another is a separate difficulty. \\
\bottomrule
\end{tabularx}
\end{table}

\section{RQ1: Challenges across modelling activities}
\label{sec:findings_rq1}

RQ1 asked what challenges students and instructors perceive, and how those challenges show up when people interpret, create, refine, and evaluate models. Four categories answer this question first: getting started (CAT1), too many symbols and frameworks (CAT2), judging whether a model is good enough (CAT12), and moving from one kind of diagram to another (CAT14). These are different problems. Waiting for a worked example is not the same as getting lost in arrows. A class diagram can look fine as UML and still drop the goals or the process that the previous model was meant to carry.

\subsection{CAT1 --- Starting is gated without a worked path}
\label{sec:cat1}

Getting from a written problem to the first real application was often gated on a demonstration, prior contact with the same notation, or some other worked path. Lecturers described a room that waits: ``They wait for the teacher to start the solution\ldots Some do, but they are a minority. The others wait'' (L1). Students described the same gate from the other side when they had not yet seen the notation used on a live task.

Wave~2 added a different form of the same gate. In S8's bachelor's project, UML and requirements models were required when S8 did not know how to use those notations, even though earlier courses had already used entity-relationship and class diagrams: ``I had the idea that I had to use UML, requirements engineering, and so on, because it was in the project requirements, but I had no idea how to use those notations'' (S8). A later Software Engineering course introduced modelling more explicitly, but by then a homemade sketching routine was already in place (CAT4).

The gate is not universal. S3 could start MSP exercises immediately because of earlier, even if shallow, BPMN contact: ``In MSP, I remember clearly that I could start the exercises as soon as the professor showed them. I did not need to wait'' (S3). S5 rarely waited for a teacher demonstration (``I rarely wait for a teacher demo''). S7 and S10 started without copying the board; wait-to-start can describe the \emph{room} (classmates who copy) without describing every student. S9, who recalled modelling at a distance, likewise needed to see the task done once before starting.

\subsection{CAT2 --- Notation and symbol overload create doubt}
\label{sec:cat2}

Unfamiliar symbols, arrows, and frameworks created doubt at the start \emph{and} mid-work. This is distinct from CAT1: S1 could name the high-level task and still stall on notation. ``The high-level concepts were fairly straightforward, but the symbols and notation complicated things'' (S1). On a game-modelling project, ``Different arrows and boxes meant different things. There was a lot of information to absorb, and it felt arbitrary when trying to do something useful with it'' (S1). S2 located course difficulty in having to learn frameworks before drawing: ``I had to use frameworks with which I was unfamiliar. I had to learn them before I could start drawing'' (S2).

Wave~2 added that similar-looking diagrams confuse until an example sits beside them (S8). When S7 and S10 got stuck moving from one kind of diagram to another, that is discussed under CAT14 rather than as more symbol overload.

\subsection{CAT12 --- Hard to judge good enough}
\label{sec:cat12}

Unlike code, a model does not compile. Students and lecturers described quality as walkthrough, justification, granularity, and readability, in a setting where several solutions can be valid. S4 used a homemade checklist because none was official: ``I have a kind of personal checklist, although nobody gave it to us as an official rubric'' (S4). S7 judged exams by a homemade list and ``faith'' that the teacher's model would match. S10 added explainability for oral defence: ``If three people who made the model had difficulty explaining it, it was probably too complicated'' (S10).

Lecturers accepted multiple valid models (L1) and still had to grade completeness, consistency, and readability. The missing compile result is what makes CAT5 and CAT6 consequential: the teacher (or a late comment) is often the only ``right/wrong'' signal (PH2).

\subsection{CAT14 --- Cross-notation switching is hard}
\label{sec:cat14}

Moving across notation families lost traceability or required a new mindset and tool set. S6's hardest moment was not drawing actors but translating an iStar rationale model into UML: ``our pair produced a class diagram that looked syntactically correct but lost the goals'' (S6), and ``nobody had explicitly explained how much of the iStar model should survive in UML'' (S6). Each family demanded a different stance: ``UML felt more structural. iStar felt organisational and political: who depends on whom, which goals conflict'' (S6).

Wave~2 densified CAT14 beyond that single RE course. S10 translated an external BPMN pool too directly into UML and was told afterwards that ``the models were answering different questions'' (S10). Names and requirements carried; process structure did not become classes. S7 reported a lighter variant: confusion at the switch was mainly syntax and duplicated teaching, not a first start. S8 got stuck when two notations looked similar (requirements diagrams versus class diagrams). That is part similar-symbol confusion (CAT2) and part moving between kinds of diagram (CAT14). Whether lecturers teach ``different questions'' as syllabus content remains open (Chapter~\ref{cha:discussion}).

\subsection{PH1 --- Turning a written problem into the first diagram}
\label{sec:ph1}

The hard, under-taught step is making the first model from a problem text---not memorising symbol names. CAT1 (missing path), CAT3 (examples as unlock), and CAT4 (self-invented procedures when teaching is thin) jointly support this phenomenon. CAT2 can delay the same first move, but it is a different mechanism: too many symbols, not the absence of a path.

\section{RQ2: Strategies and support needs}
\label{sec:findings_rq2}

RQ2 asked what people do about the challenges and what pedagogical or tool support they identify. The primary categories are teacher examples (CAT3), student-invented process tactics (CAT4), late feedback (CAT5), access to feedback (CAT6), and AI used as an explainer (CAT13).

\subsection{CAT3 --- Teacher examples and walkthroughs unlock work}
\label{sec:cat3}

Worked examples and whiteboard walkthroughs were the main pedagogical unlock. S1 put it directly: ``Examples are more useful for me because I learn by seeing things'' (S1). L1 described a board ritual: map the statement into a model step by step, then give another exercise to do alone. Wave~2 instantiated the same pattern: S7 recalled a simple board example followed by class debate; S8 said examples taught notation more than the expected answer.

A visible negative case is L4, who does not solve exercises up front: ``No, I do not solve exercises for them. They solve the exercises with me'' (L4). That practice is a pedagogical choice, not evidence that examples are unused elsewhere.

\subsection{CAT4 --- Students invent process strategies when teaching is thin}
\label{sec:cat4}

When a worked path was missing or late, students invented procedures: trial-and-error text-to-model, objects then relations, reading symbols as words, early drafts. S5's group scripted the first ten minutes before opening a tool: ``Someone reads the text aloud, we each say what jumps out, and we decide who drafts the first diagram'' (S5).

Wave~2 added tactics tied to classroom risk and project timing. S7 attacked the problem quickly ``to have something done in case I was the `victim' of the board'' (S7), while classmates copied. S8 underlined requirements and moved from pencil to pen on the bachelor's project that required UML they did not yet know how to use; earlier courses had already used entity-relationship and class diagrams (CAT1). S10 drafted in a notebook and only then put work on the board. These are student strategies, not evidence that the course designed them.

\subsection{CAT5 --- Feedback arrives too late}
\label{sec:cat5}

End-of-course or end-delivery feedback let errors compound. S1 wanted problems caught earlier: ``Problems could have been addressed along the way instead of propagating until the end'' (S1). Lecturers often planned several rounds of feedback, but class time and ECTS then cut those rounds short.

What students could still change depended on remaining time and whether the grade was already locked. S8's individual project received walk-around comments that ``almost always implied changes, large or small'' (S8), with no remembered full redo. S10 still had time before the next delivery and did ``a large local restructuring'' (S10). S7 received the grade before comments---``The grade was given before the feedback'' (S7)---and a rushed teammate then forced a full rebuild: ``we had to completely redo the system from the first phase'' (S7). CAT9 (group tempo) interacts with CAT5 here.

\subsection{CAT6 --- Feedback access depends on initiative and format}
\label{sec:cat6}

Detailed feedback often arrived only if students asked, and public formats favoured those willing to expose a half-finished diagram. S4 preferred private channels: ``Not very comfortable in public. I am shy, so presenting a half-finished diagram to the whole class is stressful'' (S4). S5, more outgoing, used a public lab board. Wave~2 widened the conditions beyond a shy/outgoing split: S7 was comfortable showing \emph{models} (more than code) but still treated extra comments as asked-for and in person; S8 shared easily with friends and the teacher, not with course-group strangers; S9 would hide poor drawing rather than a wrong model; S10 found desk walk-around easy and raising a hand hard: ``If the teacher comes by your desk and asks `is everything all right?', it is much easier to show what you are doing than to raise your hand in front of the class'' (S10). Graded versus ungraded work also gated whether S8 sought checks.

S10 wanted a private way to check a model during the lab. Whether that would have answered the small questions that never reached the teacher is a suggestion for teaching design, not a finding (S10).

\subsection{CAT13 --- AI used as explainer, not generator}
\label{sec:cat13}

Students who used AI asked it to explain options or compare approaches; they rejected generation for graded work. S5 used a chatbot ``to suggest alternative class diagrams from a bullet list, then compared options with the team'' (S5) and would not submit raw output. S10's incident is the first Portuguese lived case in the corpus: screenshots and small questions after a self-check, never a full model---``I never asked `make me the model of this whole system'.'' (S10). That use also substituted for some public questions S10 avoided (CAT6). S4 and S6 reported the same explain-not-generate stance. One lecturer wanted an LLM that would explain a student's model back in ordinary language, so the student could compare that explanation with what they had meant---not a tool that draws the assignment (L3). That is a teaching idea, not a CAT13 finding; Chapter~\ref{cha:discussion} returns to it.

\subsection{PH2 --- The teacher is the only compile button}
\label{sec:ph2}

Because a model does not compile (CAT12), students often learn they are wrong when a teacher---or late feedback---tells them (CAT5, CAT6). Walk-around, private messages, and public boards change \emph{who} gets that signal in time, not the fact that the signal is social.

\subsection{PH3 --- Want checks and explanations, not auto-drawn homework}
\label{sec:ph3}

Across tools and AI, participants asked for validation and help in understanding, not a machine that draws the assignment. CAT7 (semantic checking valued), CAT8 (friction when tools are heavy), and CAT13 (AI as explainer) feed this phenomenon. It is taken up again under RQ3 for the tool side.

\section{RQ3: Socio-technical conditions}
\label{sec:findings_rq3}

RQ3 asked how tools, course design, domains, and feedback practices shape modelling experience. The categories here are semantic checking (CAT7), tool friction (CAT8), groups (CAT9), real audiences (CAT10), and course load (CAT11). Course load cuts time for demonstrations and feedback.

\subsection{CAT7 --- Semantic checking is valued}
\label{sec:cat7}

Tools and checks that catch invalid models were wanted. This is not a claim that heavy tools always beat drawing. S3 valued a BPMN-specific interface: ``it presents a lot of the vocabulary directly in the interface'' (S3). Lecturers valued checkers such as piStar for iStar (L4) and preferred modelling tools over generic drawing when licensing allowed (L1).

Wave~2 added bounds. S7 saw an open-source compile tool presented and then dropped it: ``the interface was a bit antiquated and using it created more problems than it solved'' (S7). S8 skipped optional Camunda/Archi suites and used draw.io. Wanting checks (CAT7) and abandoning a checker for UI or optionality are compatible; they are not a new category.

\subsection{CAT8 --- Tool friction, switching, and cost}
\label{sec:cat8}

Option overload, licensing, mid-course switches, and slow or corrupt tools raised the cost of modelling. S1 found StarUML overwhelming for a simple relation: ``The tool had so many options that it became completely overwhelming to know how to make a simple entity-to-box relation'' (S1) and preferred a simpler drawing tool. Lecturers reported licensing pushing tool changes (L1). Wave~2 added an outdated compile interface that S7 then dropped, optional Camunda and Archi that S8 skipped in favour of draw.io, and UML menus S9 found outdated. Checks are still wanted (CAT7); a tool that is too complicated, optional, or slow will still be dropped.

\subsection{CAT9 --- Groups can kill or save the model}
\label{sec:cat9}

Code-first pressure or passive teammates killed diagrams; ownership, pairs, live boards, and stand-ups kept them alive. S4's group wanted to jump to implementation on a generic shop domain while S4 was still drawing the class diagram. Wave~2: a teammate who rushed to finish phase~1 forced S7's pair into a from-scratch redo (CAT5 interaction). S10's group split the work by notation---one person did BPMN, another requirements, another UML---and the diagrams started to disagree: ``Then we had two versions of the same system'' (S10). They later added a draft-and-review step so everyone could still explain the models at the defence.

\subsection{CAT10 --- Real audience makes modelling worth it}
\label{sec:cat10}

A concrete reader---supervisor, stakeholders, teaching assistant---pulled quality up more than domain excitement alone. S4 still held that ``Interesting domains make a big difference'' (S4), but wave~2 placed audience above that: course-only models ``would probably only pass through the teacher's hands'' (S7), whereas a thesis architecture for supervisor and defence became a priority. S10 similarly treated explainability as a defence criterion (CAT12). Domain interest matters; it does not replace a reader.

\subsection{CAT11 --- Course load and theory--practice gaps}
\label{sec:cat11}

Heavy continuous assessment, missed theory, slide-only catch-up, and late recognition of modelling's value shaped what teaching could actually deliver. Lecturers described workload and continuous evaluation cutting attendance and time for demos and feedback (L1, L4). Wave~2 added a student dimension that is not the same as lecturer ECTS-scarcity: S7 tried to finish modelling quickly because a heavy semester left little interest---``because modelling was not something that interested me much, I tried to get those course tasks done as fast as possible so I could focus more on the others'' (S7). S8's paid tutoring job lowered interest more than other courses. When load is this high, there is less time for the demonstrations (CAT3) and the feedback rounds (CAT5) described earlier.

\subsection{PH4 --- A diagram survives only if someone owns it and someone reads it}
\label{sec:ph4}

Group ownership (CAT9) and a real audience (CAT10) jointly decide whether modelling stays alive and careful. A well-taught example (CAT3) does not keep a diagram alive in a group that has already switched to code.

\section{Category summary}
\label{sec:findings_summary}

Table~\ref{tab:cat_rq_summary} lists each overview category with its primary research question, a one-line claim, and the memos that most strongly support it. Those identifiers point to the analytical notes; they are not extra quotations.

\begin{table}[p]
\centering
\footnotesize
\setlength{\tabcolsep}{3pt}
\caption{Overview categories, primary research question, one-line claim, and strongest memos.}
\label{tab:cat_rq_summary}
\begin{tabularx}{\textwidth}{l l >{\raggedright\arraybackslash}p{0.42\textwidth} l}
\toprule
\textbf{ID} & \textbf{RQ} & \textbf{Claim} & \textbf{Memos} \\
\midrule
CAT1 & RQ1 & First application is gated on a demo, prior contact, or a project that required notations not yet taught. & S3M1, S8M1, L1M1 \\
CAT2 & RQ1 & Symbol and framework overload create doubt; distinct from a missing path. & S1M1, S3M3, S2M3 \\
CAT12 & RQ1 & No compile result; quality is walkthrough, justification, granularity. & S6M2, S4M2, S10M5 \\
CAT14 & RQ1 & Switching families loses trace or requires a new question, not only new symbols. & S6M1, S10M1, S7M5 \\
CAT3 & RQ2 & Teacher examples and walkthroughs unlock work (L4 withholds solved examples). & S1M2, S7M1, L4M6 \\
CAT4 & RQ2 & Students invent process tactics when teaching is thin. & S5M3, S7M1, S8M1 \\
CAT5 & RQ2 & Late feedback lets errors compound; what can still change depends on time and locked grades. & S1M5, S7M3, S10M3 \\
CAT6 & RQ2 & Access depends on initiative, graded work, and public versus desk format. & S4M3, S8M2, S10M2 \\
CAT13 & RQ2 & AI is used to explain or compare, not to generate graded models. & S4M5, S5M4, S10M4 \\
CAT7 & RQ3 & Semantic checking is valued; a heavy or optional checker may still be dropped. & S3M4, S7M2, L4M5 \\
CAT8 & RQ3 & Licensing, option load, switching, and weak UI raise modelling cost. & S1M4, S7M2, L1M5 \\
CAT9 & RQ3 & Groups kill diagrams under code-first pressure or save them via ownership. & S4M1, S7M3, S10M5 \\
CAT10 & RQ3 & A concrete reader pulls quality up more than domain excitement alone. & S7M5, S10M5, L4M3 \\
CAT11 & RQ3 & Load and theory--practice gaps block demos and feedback cycles. & L1M3, S7M6, S8M3 \\
\bottomrule
\end{tabularx}
\end{table}

The table is the findings of this chapter. Teaching ideas and open questions that did not earn a category are taken up in Chapter~\ref{cha:discussion} (Table~\ref{tab:open_questions}), not as extra rows here.
